\chapter{Conclusions and Future Work}
\section{Conclusion}

This dissertation tested whether conventional anti-spam tools work over a
delay-tolerant network. A testbed was implemented using ION-DTN, BPMail gateway, Postfix/Dovecot, local unbound DNS resolver, and \texttt{tc netem} to add delay. Three different tools – SpamAssassin, rspamd, and Vipul's Razor were deployed and evaluated against an earth-to-moon delay scenario with a non-delayed setup as the baseline. The evaluation shows that SpamAssassin, rspamd, and Vipul's Razor are designed for terrestrial network latencies. When round-trip time exceeds their expected timeout values, spam detection becomes less effective or may fail.

However, spam detection is possible with appropriate configuration of unbound. By default, unbound will fail silently without any indication: spam filter will still check the messages and return a score but will not detect spam at all. This behaviour is difficult to detect because spam filter still returns a score while DNS-based detections are silently absent. If a tool's timeout is shorter than the round-trip time, DNS-based detections may be missed rather than delayed. Blocklist matches are replaced by timeout results, while a tool that depends on a single collaborative server may return only network errors.

Delay also increases scan time. A scan waits for network lookups to complete, so a process that normally finishes in under a second can take several seconds. The cost depends on the depth of dependent lookups rather than their number: messages with long dependency chains wait for multiple round trips. This increases message latency and can cause mail queues to build up while scans wait for delayed network responses.

Network traffic also changes under delay. In SpamAssassin, DNS traffic is much larger than the email traffic itself, so DNS lookups account for most bandwidth usage. But, once the resolver stops sending upstream queries, it reduced DNS traffic but prevented DNSBL lookups.  rspamd increases DNS traffic because timed-out queries are retransmitted before delayed replies arrive. Razor's traffic is dominated by repeated TCP connection attempts rather than DNS. In all three tools, a significant proportion of network traffic is spent on lookup or connection attempts that do not contribute to spam detection.


These results were measured using an Earth--Moon delay. Earth--Mars delay is measured in minutes each way rather than seconds, so the round-trip time is far greater than the default timeout values used by all three tools. The timeout-related behaviour observed in the reduced-timeout experiments would therefore be expected under the default configurations. Messages with dependency chains, which took several seconds to scan in the Earth--Moon scenario, would require many minutes. Retransmissions would also increase DNS traffic on a link with limited bandwidth. Synchronous DNS-based spam filtering is therefore unlikely to remain practical at interplanetary distances.


At Earth--Moon distance, increasing timeout values above the round-trip time restores DNS-based spam detection, but each message scan takes several seconds and mail throughput is reduced because scans wait for delayed network responses. Beyond Earth--Moon distances, simply increasing timeout values is unlikely to remain practical. At interplanetary distances, the required timeout values would increase from seconds to many minutes, making synchronous network lookups impractical. Future spam-filtering systems for DTN may instead need to replace synchronous network lookups with asynchronous mechanisms that better match DTN's store-and-forward operation.

\section{Limitations}

This study met its objectives, but has the following limitations:

\begin{itemize}

\item \textbf{Single fixed delay.} All measurements were performed using a single
Earth--Moon round-trip delay. Larger delays, as well as delay with jitter,
packet loss, or link interruptions, were not evaluated. The discussion of
Earth--Mars distances is therefore based on the observed mechanisms rather than
direct measurement.

\item \textbf{Blocklist drift and interval between runs.} Since blocklists change
over time, comparing individual message verdicts between runs performed on
different days reflects both delay effects and changes in the blocklists.
Runs~2 and~3 were performed two days apart, so per-message verdict changes are
affected by this drift. The overall detection totals for each run are affected
much less because blocklist drift changes only a small proportion of messages.

\item \textbf{Delay only, not a complete deep-space network.} Delay was introduced
using \texttt{tc netem} in a one-node ION-DTN testbed. The testbed accurately
emulates high round-trip time but does not reproduce other characteristics of
deep-space communication, such as scheduled communication windows, prolonged
link interruptions, or higher transmission error rates. The findings therefore
reflect the effects of propagation delay rather than all characteristics of a
deep-space network.

\item \textbf{Version-dependent behaviour.} The results apply to the specific
software versions used in this study. For example, unbound~1.13.1 does not
provide an \texttt{infra-cache-max-rtt} setting, so the resolver behaviour
observed in Run~3 could not be limited through configuration. Later versions
introduce this option, so behaviour may differ.
\end{itemize}


\section{Future Work}

The findings and limitations of this study suggest the following directions for future work:

\begin{itemize}

\item \textbf{Evaluate Mars-scale delays and more challenging network conditions.}
The conclusions for Earth--Mars communication are based on the mechanisms
observed in this study rather than direct measurement. Future work should
evaluate the tools under minute-scale round-trip times together with jitter,
packet loss, and link interruptions to determine whether the observed behaviour
extends to these conditions.

\item \textbf{Move network-dependent filtering off the per-message path.} The
tools evaluated in this study perform live network lookups for each message,
which leads to the latency and bandwidth costs observed in this work. One
approach is to maintain a local copy of blocklist data at the receiving site,
allowing per-message checks to be performed locally while synchronising the
database asynchronously over the DTN path. Another approach is to perform the
checks on the Earth side, where blocklist servers are directly reachable, and
attach the results to the message before transmission. The first approach avoids relying on a trusted upstream but uses blocklist data that may be out of date. The second uses current blocklist data. The second uses current blocklist data, but it requires the receiving organisation to trust another organisation's spam-filtering decisions before accepting the message.

\item \textbf{Deferred asynchronous filtering.} Instead of delaying delivery while waiting for network lookups, a message could be accepted and stored immediately, with DNS or collaborative-server lookups performed asynchronously afterwards. The spam classification could then be updated once the lookups complete, for example by moving the message to a quarantine folder or marking it as spam. This approach removes the delay from the delivery path, although spam classification is no longer available immediately.

\end{itemize}